\chapter{Linear Momentum}

\mathTopics{\pgRef{chap:vectors}}

\section{Momentum}

\textbf{Momentum}~(symbol: \(\vecb{p}\)) has SI units of \unit{\kilogram
\metre \per \second}
or \unit{\newton \second},
and is given by the formula
\begin{equation}
    \vecb{p} = m \vecb{v},
\end{equation}
where \(m\) is the mass and \(\vecb{v}\) is the velocity.

A system's momentum is the sum of its parts' momenta, i.e.,
\begin{equation}
    \vecb{p}_\mathrm{system} = \sum_i m_i \vecb{v}_i.
\end{equation}

For objects with constant mass,
\begin{equation}
    \Delta{\vecb{p}} = m \Delta{\vecb{v}}.
\end{equation}

\begin{equation} \label{eq:force-momentum}
    \begin{aligned}
        \odv{\vecb{p}}{t} &= m \odv{\vecb{v}}{t} \\
        \odv{\vecb{p}}{t} &= m \vecb{a} \\
        \vecb{F} &= \odv{\vecb{p}}{t}
    \end{aligned}
\end{equation}

In integral form,
\begin{equation}
    \int \vecb{F}_\mathrm{net} \odif{t} = \Delta{\vecb{p}}.
\end{equation}

\section{Impulse}

\textbf{Impulse}~(symbol: \(\vecb{J}\)) has SI units of \unit{\newton \second}
or \unit{\kilogram \metre \per \second}.
In the simple case where \(\odv{\vecb{F}}{t} = 0\),
\begin{equation}
    \vecb{J} = \vecb{F} t.
\end{equation}
The more general equation is
\begin{equation}
    \vecb{J} = \int \vecb{F}(t) \odif{t} = \Delta{\vecb{p}},
\end{equation}
where \(\Delta{\vecb{p}}\) is the change in momentum.

\section{Conservation of Momentum}

Because of \pgRef{eq:force-momentum}
and \hyperref[itm:newtons-1st-law]{Newton's first law of motion}
(\cpageref{itm:newtons-1st-law}),
when there is no external force, \(\Delta{\vecb{p}} = \vecb{0}\),
i.e, the final and initial momentum have to be equal.

Types of \textbf{collisions}:
\begin{itemize}
    \item
        \textbf{Elastic}
        \begin{itemize}
            \item Kinetic energy is conserved.
            \item Objects bounce apart with the same energy they collided with.
        \end{itemize}
    \item
        \textbf{Inelastic}
        \begin{itemize}
            \item Kinetic energy is lost.
            \item Objects bounce apart with less energy or not all.
                The case where the objects do not bounce at all
                is known as \textbf{perfectly inelastic}
        \end{itemize}
    \item
        \textbf{Explosion}
        \begin{itemize}
            \item Kinetic energy is gained.
            \item Objects bounce apart.
        \end{itemize}
\end{itemize}

Practically speaking, no collision is perfectly elastic or inelastic;
they are on a sliding continuum between the two.
